%
% This work is licensed under a Creative Commons Attribution-ShareAlike 4.0 International License.
%

\documentclass[10pt, letterpaper, twoside, openany]{article}

\title{Introduction to Quantum Mechanics -- Worksheet 5}
\author{Joseph D. MacMillan}
\date{}


\usepackage{graphicx}
\usepackage{color} 
\usepackage{amsmath, amssymb}
\usepackage{libertinus}
\usepackage{microtype}
\usepackage{layout}
\usepackage[most]{tcolorbox}
\tcbuselibrary{skins,breakable}
\usepackage[hang, small, bf]{caption}
\captionsetup[table]{position=top}
\captionsetup[figure]{position=bottom}
\usepackage[hyphens]{url}
\usepackage[hidelinks]{hyperref}
\hypersetup{pdftitle={Introduction to Quantum Mechanics -- Worksheet 5}}
\hypersetup{pdfauthor={Joseph D. MacMillan}}
\urlstyle{rm}
\usepackage[shortlabels]{enumitem}
\usepackage[margin=1in]{geometry}
\usepackage{fancyhdr}




\newcommand{\uu}{\symbfup{u}}
\newcommand{\grad}{\symbfup{\nabla}}
\newcommand{\curl}{\symbfup{\nabla} \times}
\newcommand{\dfdx}[2]{\frac{\partial {#1}}{\partial {#2}}}
\newcommand{\ddfdx}[2]{\frac{\partial^2 {#1}}{\partial {#2}^2}}
\newcommand{\vort}{\symbfup{\omega}}
\renewcommand\vec{\symbfup}
\newcommand{\unit}[1]{\hat{\vec{#1}}}
\newcommand{\U}{\mathbb{U}}
\newcommand{\V}{\mathbb{V}}
\newcommand{\LL}{\mathbb{L}^2}
\newcommand{\LLLL}{\mathbb{L}^4}

\newcommand{\ket}[1]{| #1 \rangle}
\newcommand{\bra}[1]{\langle #1 |}
\newcommand{\braket}[2]{\langle #1 | #2 \rangle}
\newcommand{\avg}[1]{\langle #1 \rangle}


%\DeclareMathAlphabet{\mathbfsf}{\encodingdefault}{\sfdefault}{bx}{sl}
\newcommand{\tens}[1]{\mathbfsfit{#1}}


\newcounter{example}[section]
\def\theexample{\thechapter.\arabic{example}}
\newenvironment{example}[1][ ]{\refstepcounter{example}
\begin{tcolorbox}[breakable, sharp corners, boxrule = 0pt, frame empty, opacityframe=0, parbox=false]
\textbf{Example \theexample \ -- #1.}
}
{ 
\end{tcolorbox}
}

\newenvironment{theorem}[1][ ]{
\begin{tcolorbox}[breakable, sharp corners, boxrule = 0pt, frame empty, opacityframe=0, parbox=false]
\textbf{#1.}
}
{ 
\end{tcolorbox}
}

\newcounter{problem}[section]
\def\theproblem{\thechapter.\arabic{problem}}
\newenvironment{problem}[1][ ]{\refstepcounter{problem} \noindent \textbf{Problem \theproblem \ -- #1.}}{\vspace{0.1in}}

\renewcommand{\arraystretch}{2.5}

\begin{document}

\pagestyle{fancy}
\fancyhead[L]{\emph{Introduction to Quantum Mechanics}}
\fancyhead[R]{Worksheet 5} 

\Large \noindent Worksheet 5

\noindent \LARGE \textbf{Time Evolution}

\normalsize

\vspace{0.2in}

\noindent Operators  \bullet \  Eigenvalues   \bullet \ Hamiltonian \bullet \ Schr\"odinger \bullet \ Evolving states  

\vspace{0.1in}

\noindent \rule{\textwidth}{0.5pt}

\vspace{0.2in}

\begin{enumerate}

\item The Hamiltonian of a two-state system is given by
\[
\hat{H} \to \begin{pmatrix}
E_1 & 0 \\ 0 & E_2
\end{pmatrix}.
\]

What basis is this Hamiltonian written in?  How do you know?  

\vspace{1in}

\item What do the values $E_1$ and $E_2$ represent?  What do we call them?  Write the energy eigenstates -- call them $\ket{E_1}$ and $\ket{E_2}$ -- in matrix notation.

\vspace{1in}

\item A second observable, which we'll call $A$, can be written in the energy basis as
\[
\hat{A} \to \begin{pmatrix}
0 & a \\ a & 0
\end{pmatrix}.
\]
Do $\hat{H}$ and $\hat{A}$ commute?  What does that imply about the two observables?

\newpage

\item Find the eigenvalues and eigenstates of the operator $\hat{A}$.  What do the eigenvalues represent?  What about the eigenstates?

\vspace{5in}

\item Suppose at time $t=0$, our quantum system is in a state given by
\[
\ket{\psi(0)} = \ket{+a}.
\]
If we measure $A$ on this state, what are the possible results and the probabilities of each?

\newpage

\item Write the state $\ket{\psi(0)}$ in the energy basis, both in Dirac and matrix notation.  What is the exact value of the energy of the system at $t=0$?

\vspace{1in}

\item What if we go ahead and measure the energy $E$ -- what are the possible results and the probabilities of each?

\vspace{2in}

\item What is the state at a later time?  Write it in the energy basis.

\vspace{1in}

\item Write the state using the eigenstates of $A$ as your basis.

\newpage

\item At time $t$, a measurement of energy is made.  What are the possible results and probabilities?  Comment on them.

\vspace{1in}

\item At time $t$, a measurement of the observable $A$ is made.  What is the probability of measuring $A = +a$?  

\end{enumerate}

\end{document}
